\documentclass{beamer}
\usetheme{Madrid}
\usecolortheme{whale}

\usepackage[utf8]{inputenc}
\usepackage[spanish]{babel}
\usepackage{amsmath}
\usepackage{amsfonts}
\usepackage{amssymb}

% --- Título y Datos ---
\title{George Boole}
\subtitle{La lógica como práctica y herramienta social}
\author{Karina G. Buendía}
\date{abril de 2026}

\begin{document}

% Slide 1: Portada
\frame{\titlepage}

% Slide 2: Introducción - ¿Qué es la matemática situada?
\begin{frame}{Matemática basada en prácticas}
    \begin{itemize}
        \item \textbf{Perspectiva:} La matemática no surge en el vacío, sino de necesidades humanas y contextos específicos.
        \item \textbf{Boole como practicante:} Su obra no fue solo abstracción pura, sino un intento de democratizar el pensamiento lógico.
        \item \textbf{El objetivo:} Reducir el razonamiento a un "cálculo" manipulable, similar al álgebra tradicional.
    \end{itemize}
\end{frame}

% Slide 3: Contexto Social - El autodidacta
\begin{frame}{El origen: Lincoln y la lucha de clases}
    \begin{block}{Un matemático fuera de la élite}
        A diferencia de sus contemporáneos en Cambridge, Boole fue un \textbf{autodidacta} de clase trabajadora.
    \end{block}
    \begin{itemize}
        \item Hijo de un zapatero con intereses científicos.
        \item La imposibilidad económica de asistir a la universidad marcó su "práctica": la necesidad de aprender y enseñar fuera de los muros académicos.
        \item \textbf{Práctica situacional:} Fundó su propia escuela a los 19 años para sostener a su familia.
    \end{itemize}
\end{frame}

% Slide 4: El giro hacia la Lógica
\begin{frame}{La Lógica como Álgebra}
    Boole observó que las leyes del pensamiento seguían patrones similares a las operaciones numéricas.
    \vspace{0.3cm}
    \begin{columns}
        \begin{column}{0.5\textwidth}
            \textbf{Práctica Algebraica:}
            Utilizó símbolos para representar conceptos (clases) y operadores para las relaciones.
        \end{column}
        \begin{column}{0.5\textwidth}
            \begin{exampleblock}{Leyes Fundamentales}
                \begin{itemize}
                    \item $x \cdot y = y \cdot x$ (Conmutativa)
                    \item $x^2 = x$ (Ley de Idempotencia)
                \end{itemize}
            \end{exampleblock}
        \end{column}
    \end{columns}
\end{frame}

\begin{frame}{Investigación sobre las leyes del pensamiento}
    En su obra de 1854, Boole formaliza la idea de que la lógica es una rama de la matemática.
    \begin{itemize}
        \item \textbf{El Universo ($1$):} Todo lo existente.
        \item \textbf{La Nada ($0$):} El conjunto vacío.
        \item \textbf{Selección:} Las variables $x, y, z$ actúan como operadores que seleccionan objetos de un grupo.
    \end{itemize}
\begin{block}{Cita}
“Las leyes de los símbolos son las leyes de las operaciones que ellos representan”.
\end{block}

\end{frame}

% Slide 6: Boole en Irlanda - Queen's College Cork
\begin{frame}{La institucionalización de la práctica}
    En 1849, Boole es nombrado profesor en Cork, Irlanda. 
    \begin{itemize}
        \item \textbf{Contexto:} Irlanda post-hambruna. Su práctica docente se enfocó en la rigurosidad y la accesibilidad.
        \item \textbf{Interacción:} Su correspondencia con Augustus De Morgan influyó en la evolución de su sistema lógico.
        \item Su muerte prematura en 1864 (causada por una neumonía tras caminar bajo la lluvia para dar clase) refleja su entrega a la labor pedagógica.
    \end{itemize}
\end{frame}

% Slide 7: Legado - De la filosofía al silicio
\begin{frame}{Legado: La práctica que construyó el futuro}
    Aunque Boole veía su obra como una exploración de la mente humana, su "práctica" matemática permitió:
    \begin{enumerate}
        \item \textbf{Claude Shannon (1937):} La aplicación del álgebra de Boole a los circuitos de conmutación.
        \item \textbf{Arquitectura Computacional:} Todo procesador moderno opera bajo la lógica binaria de Boole.
        \item \textbf{Búsqueda de Información:} Operadores AND, OR, NOT.
    \end{enumerate}
\end{frame}

% Slide 8: Conclusión
\begin{frame}{Conclusión}
    \begin{alertblock}{Reflexión Final}
        La biografía de Boole demuestra que la matemática basada en prácticas es una herramienta de empoderamiento. Su álgebra no fue un juego de símbolos, sino una búsqueda de las estructuras universales del razonamiento humano.
    \end{alertblock}
\end{frame}

\end{document}